\chapter{Desarrollo}

El cuerpo principal de la memoria describe el trabajo realizado.
En este ejemplo, demostramos algunas construcciones típicas: una
cita bibliográfica~\parencite{knuth1986} y referencias cruzadas a
la sección~\ref{sec:arquitectura}. La cita anterior es parentética
(autor y año dentro del paréntesis); cuando el autor ya forma parte
de la frase conviene usar \verb|\textcite|, que sólo encierra el
año: \textcite{tolkien1954lord} cuenta cómo Bilbo decía a menudo
que solo había un camino y que era como un río caudaloso; nacía
en el umbral de todas las puertas, y todos los senderos eran ríos
tributarios. "Es muy peligroso, Frodo, cruzar la puerta", solía
decirme. "Vas hacia el Camino, y si no cuidas tus pasos no sabes
hacia donde te arrastrarán".

\section{Arquitectura}
\label{sec:arquitectura}

La arquitectura del sistema sigue el patrón de separación entre
normativa y estilo. La clase \texttt{tfe-ucmfdi} implementa los
requisitos obligatorios de la normativa UCM-FDI; el aspecto
visual concreto queda delegado a la opción \texttt{estilo}.

\section{Implementación}

La implementación se basa en la clase KOMA-Script
\texttt{scrbook} \autocite{kohm2004koma}, con soporte nativo para UTF-8 mediante
LuaLaTeX.
